\section{引言}
双臂遮挡操作包含一组相互依赖的交互：移开遮挡物以发现目标，将两者分离，
把目标运输至指定区域，并恢复遮挡物。外观相近的画面可能对应不同的接触关系或任务历史；
同一个先前合理的推动动作，也可能在目标移动或重新遮挡后不再合适。
因此，任务的难点不仅是复现示范轨迹，而是在执行过程中持续判断
\emph{当前交互关系是否支持继续执行原有动作}。

动作块模仿学习通过一次预测多步动作提高运动连贯性，ACT是其中一种代表\cite{zhao2023act}。
这类方法可以周期重规划，并非天然没有反馈。然而，当计划以较低频率更新、
而中间若干动作按原计划执行时，新发生的接触变化无法立即影响这些指令。
同时，动作回归本身不要求内部表示显式区分目标可见程度、目标与工具的相对位置，
或任务是否已接近完成。这使得拨空气、错过接触后继续推动和完成附近的无效运动
成为需要分析的行为，而非仅以平均动作loss描述的误差。

本文提出的研究问题是：\emph{在有限示范与较少额外人工标注下，如何利用可得观测与
交互历史，学习与真实任务状态相适应的动作，提高遮挡接触操作的完成能力？}
输入图像、任务相关token和历史条件残差是达到这一共同目标的不同途径，
而不是三个互不相关的任务。我们将其拆成两个互补缺口。
其一是有限样本下任务相关关系可能未被可靠编码的\emph{表示缺口}；
其二是计划观测已经陈旧、当前单帧又可能无法辨别运动趋势的\emph{反馈缺口}。
两个缺口不能互相替代：准确分割不会自行更新旧动作，增加反馈频率也不能保证
模型学会辨别错误接触。

已有研究表明，模型能否解释监督输入本身需要检验\cite{bobu2020}；
多模态模仿学习也需要区分数据构建、观测模态与执行机制各自的作用\cite{ablett2025}。
本文据此将标签可靠性与在线反馈分别建模，不把增加输入模态作为充分论据。

我们的分析以同一动作预测风险为基础，分别给出三种改进条件。
新增RGB的条件均值信息收益必须超过编码和学习代价；
同帧语义图虽不创造新信息，但通过几何监督可以改善任务相关表示，
当其表示与读出误差上界低于基线的信息损失下界时，获得严格比较保证；
历史条件残差则在可预测误差能量超过记忆编码与拟合损失时改善基础预测。
这些结论解释何种条件足以支持改善，不把实机观察到的提升直接改写为无条件定理。

据此，方法包含一条连续的信息链：通过自主提示、SAM2传播、质量诊断及局部迭代
构建标签；训练冻结U-Net提供在线语义；以几何监督查询token将任务关系写入ACT表示；
最后冻结慢策略，以最新视觉、慢语义上下文和因果历史训练有界GRU残差。
这里的进展表示是可测的几何线索，不等同于完整阶段状态机或独立完成检测器。
任务级进展能力应由事件附近的行为、可读出准确率及对应干预实验共同支持。

本文围绕三个研究要点组织论证：第一，将遮挡接触操作中的统一学习目标形式化，
给出观测、语义表示与历史残差的条件性比较保证；第二，构建从低人工成本语义监督
到几何表示，再到时间对齐历史残差的模块化方法；第三，分别检验标签可靠性、
几何利用、新观测与历史贡献、以及时延下的纠偏效果，并以跨平台任务结果评估其实际价值。
标准条件期望恒等式、一般残差结构和GRU单元本身不作为独创贡献。
已有实机测试支持语义输入、任务相关token及GRU残差设计的效果，
本文进一步用理论条件解释收益来源；当前结果不能自动证明所有条件在任意任务成立。
